\begin{description}
\item[(1)] Using the virial theorem for our isolated, spherical system of $N$
galaxies of mass $m$, each, we get
\[
-2\langle K\rangle = \langle U\rangle
\;\Rightarrow\;
-\frac{2}{N}\sum_{1}^{N}\frac{1}{2}m_i u_i^2 = \;<U>
\;\Rightarrow\;
-\frac{m}{N}\sum_{1}^{N}u_i^2 = \frac{U}{N}
\]
\hfill(3 point)

where
$\dfrac{1}{N}\displaystyle\sum_{1}^{N}u_i^2 \approx \langle u^2\rangle
= \langle u_r^{\,2}\rangle + \langle u_\theta^{\,2}\rangle
+ \langle u_\phi^{\,2}\rangle$,
where $u_r$, $u_\theta$ and $u_\phi$ are the radial velocity and the two
perpendicular velocities on the plane of the sky of the members of the
cluster.
\hfill(6 points)

Assuming that
$\langle u_r^{\,2}\rangle \sim \langle u_\theta^{\,2}\rangle
\sim \langle u_\phi^{\,2}\rangle$, we have
\[
-3m\langle u_r^2\rangle = U/N
= \Bigl(-\frac{3}{5}\frac{GM^2}{R}\Bigr)/N
= -\frac{3}{5}\frac{GMm}{R}
\;\Rightarrow\;
M \approx \frac{5R\langle u_r^2\rangle}{G}
\]
\hfill(13 points)

(where we used the gravitational potential of a spherical homogeneous mass
$M$ enclosed within radius $R$)

Alternatively the student can give a rougher order of magnitude estimate
\[
-2\langle K\rangle = \langle U\rangle
\;\Rightarrow\;
-2.\frac{1}{2}\langle u_r^2\rangle \approx -\frac{GM}{R}
\quad\text{etc}
\]
If they do not use the exact formula:
\hfill(11 Points)

\item[(2)] From the result of the previous question we have
$M \approx \dfrac{5R\langle u_r^2\rangle}{G}$, where
\[
\sqrt{\langle u_r^2\rangle} = \sigma = 1000\ km\,s^{-1}
\]
\hfill(8 points)

The angular diameter of the cluster is $\varphi = 4^\circ$ at a distance of
$d = 90$ Mpc. Therefore the diameter D of the cluster in Mpc is calculated
from:
\[
\tan\phi = \frac{D}{d}
\;\Rightarrow\;
\phi(rad) \approx \frac{D}{d}
\;\Rightarrow\;
D \approx 4\times\frac{\pi}{180}\times90\,\mathrm{Mpc} \approx 6.3\,Mpc
\;\Rightarrow\;
R \approx 3\,Mpc
\]
\hfill(11 points)

Therefore
$M \approx \dfrac{5R\langle u_r^2\rangle}{G} = 3.6\times10^{15}\,\mathrm{M}_\odot$
\hfill(3 point)

\item[(3)]
\[
\frac{M_{virial}}{L_{galaxies}}
= \frac{3.6\times10^{15}\,\mathrm{M}_\odot}{5\times10^{12}\,L_\odot}
= 720\,\frac{\mathrm{M}_\odot}{L_\odot}
\]
This is obviously much larger from
$\dfrac{\mathrm{M}_\odot}{L_\odot} = 1$ which is found from the visible mass
of the cluster. Thus $\dfrac{M}{M_{virial}} = \dfrac{1}{720}$.
\hfill(6 points)
\end{description}
